\section*{Solution}

Consider the bimatrix game
\[
\begin{array}{c|ccc}
 & L & M & R\\
\hline
U & (4,1) & (3,6) & (a,7)\\
D & (5,6) & (2,5) & (5,3)
\end{array}
\]
where the row player chooses \(U,D\) and the column player chooses \(L,M,R\).

We solve the game for \(a=3\), \(a=6\), and \(a=5\).

\subsection*{Step 1: Best responses in general}

For the row player:

\begin{itemize}
\item Against \(L\): \(4<5\), so the best response is \(D\)
\item Against \(M\): \(3>2\), so the best response is \(U\)
\item Against \(R\): compare \(a\) and \(5\):
\[
\begin{cases}
D, & a<5,\\
U \text{ and } D, & a=5,\\
U, & a>5
\end{cases}
\]
\end{itemize}

For the column player:

\begin{itemize}
\item Against \(U\): compare \(1,6,7\), so the best response is \(R\)
\item Against \(D\): compare \(6,5,3\), so the best response is \(L\)
\end{itemize}

Thus:
\begin{itemize}
\item \((D,L)\) is always a pure Nash equilibrium
\item \((U,R)\) is a pure Nash equilibrium if \(a\ge 5\)
\end{itemize}

\subsection*{Step 2: Mixed equilibrium conditions}

Let the row player choose \(U\) with probability \(p\), and the column player choose
\[
L,M,R
\]
with probabilities
\[
x,y,z,
\qquad x+y+z=1
\]

For the row player to mix, he must be indifferent between \(U\) and \(D\):
\[
4x+3y+az = 5x+2y+5z
\]
Hence
\[
y = x + (5-a)z \tag{1}
\]

For the column player, expected payoffs are:
\[
u_C(L)=p\cdot 1+(1-p)\cdot 6=6-5p,
\]
\[
u_C(M)=p\cdot 6+(1-p)\cdot 5=5+p,
\]
\[
u_C(R)=p\cdot 7+(1-p)\cdot 3=3+4p
\]

Equality conditions:
\[
u_C(L)=u_C(M) \iff 6-5p=5+p \iff p=\frac16,
\]
\[
u_C(M)=u_C(R) \iff 5+p=3+4p \iff p=\frac23,
\]
\[
u_C(L)=u_C(R) \iff 6-5p=3+4p \iff p=\frac13
\]

We now analyze each case separately.

\section*{(a) Case \(a=3\)}

The matrix is
\[
\begin{array}{c|ccc}
 & L & M & R\\
\hline
U & (4,1) & (3,6) & (3,7)\\
D & (5,6) & (2,5) & (5,3)
\end{array}
\]

\subsection*{Pure Nash equilibria}

Since \(a=3<5\), against \(R\) the row player prefers \(D\). Therefore \((U,R)\) is not a Nash equilibrium.

So the only pure Nash equilibrium is
\[
(D,L).
\]

\subsection*{Mixed Nash equilibria}

For \(a=3\), equation (1) becomes
\[
y=x+2z. \tag{2}
\]

\paragraph{Mixing between \(L\) and \(M\).}
This requires \(p=\frac16\), and then \(L\) and \(M\) are optimal for the column player.

Set \(z=0\). Then from (2),
\[
y=x.
\]
Together with \(x+y=1\), we get
\[
x=y=\frac12.
\]

Hence one mixed equilibrium is:
\[
p=\frac16,\qquad (x,y,z)=\left(\frac12,\frac12,0\right).
\]

\paragraph{Mixing between \(M\) and \(R\).}
This requires \(p=\frac23\), and then \(M\) and \(R\) are optimal for the column player.

Set \(x=0\). Then from (2),
\[
y=2z.
\]
Together with \(y+z=1\), we get
\[
z=\frac13,\qquad y=\frac23.
\]

Hence another mixed equilibrium is:
\[
p=\frac23,\qquad (x,y,z)=\left(0,\frac23,\frac13\right).
\]

\paragraph{Mixing between \(L\) and \(R\).}
This would require \(p=\frac13\). But then
\[
u_C(M)=5+\frac13=\frac{16}{3},
\qquad
u_C(L)=u_C(R)=\frac{13}{3},
\]
so \(M\) is strictly better. Thus no equilibrium of this type exists.

\subsection*{Answer for \(a=3\)}

Pure Nash equilibrium:
\[
(D,L).
\]

Mixed Nash equilibria:
\[
\left(\frac16 U+\frac56 D,\ \frac12 L+\frac12 M\right),
\]
\[
\left(\frac23 U+\frac13 D,\ \frac23 M+\frac13 R\right).
\]

\section*{(b) Case \(a=6\)}

The matrix is
\[
\begin{array}{c|ccc}
 & L & M & R\\
\hline
U & (4,1) & (3,6) & (6,7)\\
D & (5,6) & (2,5) & (5,3)
\end{array}
\]

\subsection*{Pure Nash equilibria}

Since \(a=6>5\), against \(R\) the row player prefers \(U\). Hence \((U,R)\) is a Nash equilibrium.

Therefore the pure Nash equilibria are
\[
(D,L),\qquad (U,R).
\]

\subsection*{Mixed Nash equilibria}

For \(a=6\), equation (1) becomes
\[
y=x-z. \tag{3}
\]
Equivalently,
\[
x=y+z.
\]
Using \(x+y+z=1\), we obtain
\[
(y+z)+y+z=1
\quad\Rightarrow\quad
2y+2z=1
\quad\Rightarrow\quad
x=\frac12,\qquad y+z=\frac12.
\]

\paragraph{Mixing between \(L\) and \(M\).}
This requires \(p=\frac16\). Set \(z=0\). Then
\[
x=\frac12,\qquad y=\frac12.
\]
So we get the mixed equilibrium
\[
p=\frac16,\qquad (x,y,z)=\left(\frac12,\frac12,0\right).
\]

\paragraph{Mixing between \(M\) and \(R\).}
This would require \(p=\frac23\), but such a mixture has \(x=0\), contradicting \(x=\frac12\). So no such equilibrium exists.

\paragraph{Mixing between \(L\) and \(R\).}
This would require \(p=\frac13\), but then \(M\) is strictly better than both \(L\) and \(R\). So no such equilibrium exists.

\subsection*{Answer for \(a=6\)}

Pure Nash equilibria:
\[
(D,L),\qquad (U,R).
\]

Mixed Nash equilibrium:
\[
\left(\frac16 U+\frac56 D,\ \frac12 L+\frac12 M\right).
\]

\section*{(c) Case \(a=5\)}

The matrix is
\[
\begin{array}{c|ccc}
 & L & M & R\\
\hline
U & (4,1) & (3,6) & (5,7)\\
D & (5,6) & (2,5) & (5,3)
\end{array}
\]

\subsection*{Pure Nash equilibria}

Since \(a=5\), against \(R\) the row player is indifferent between \(U\) and \(D\). Therefore \((U,R)\) is a Nash equilibrium, and as before \((D,L)\) is also a Nash equilibrium.

Thus the pure Nash equilibria are
\[
(D,L),\qquad (U,R).
\]

\subsection*{Mixed Nash equilibria}

For \(a=5\), equation (1) becomes
\[
y=x. \tag{4}
\]

\paragraph{Mixing between \(L\) and \(M\).}
This requires \(p=\frac16\). Set \(z=0\). Then by (4),
\[
x=y.
\]
Since \(x+y=1\), we get
\[
x=y=\frac12.
\]
Hence one mixed equilibrium is
\[
p=\frac16,\qquad (x,y,z)=\left(\frac12,\frac12,0\right).
\]

\paragraph{Column player uses only \(R\).}
Suppose the column player plays \(R\) with probability \(1\), i.e.
\[
(x,y,z)=(0,0,1).
\]
Against \(R\), the row player gets
\[
u_R(U)=5,\qquad u_R(D)=5,
\]
so any mixture of \(U\) and \(D\) is a best response.

Now we need \(R\) to be a best response for the column player. This means
\[
u_C(R)\ge u_C(L),\qquad u_C(R)\ge u_C(M).
\]
That is,
\[
3+4p \ge 6-5p,
\qquad
3+4p \ge 5+p.
\]
These inequalities give
\[
9p\ge 3 \iff p\ge \frac13,
\qquad
3p\ge 2 \iff p\ge \frac23.
\]
Hence the condition is
\[
p\ge \frac23.
\]

Therefore, for every
\[
p\in\left[\frac23,1\right],
\]
the profile
\[
\left(pU+(1-p)D,\ R\right)
\]
is a Nash equilibrium.

\paragraph{Mixing between \(M\) and \(R\).}
If \(x=0\), then from (4) also \(y=0\), so \(z=1\). This is exactly the previous family with pure \(R\); there is no additional interior mixture of \(M\) and \(R\).

\paragraph{Mixing between \(L\) and \(R\).}
This would require \(p=\frac13\), but then \(M\) is strictly better than both \(L\) and \(R\), so no such equilibrium exists.

\subsection*{Answer for \(a=5\)}

Pure Nash equilibria:
\[
(D,L),\qquad (U,R).
\]

Additional mixed equilibria:
\[
\left(\frac16 U+\frac56 D,\ \frac12 L+\frac12 M\right),
\]
and the whole family
\[
\left(pU+(1-p)D,\ R\right),
\qquad
p\in\left[\frac23,1\right].
\]

\section*{Final answers}

\[
\text{For }a=3: \ (D,L),\ \left(\frac16 U+\frac56 D,\frac12 L+\frac12 M\right),\ \left(\frac23 U+\frac13 D,\frac23 M+\frac13 R\right)
\]

\[
\text{For }a=6: \ (D,L),\ (U,R),\ \left(\frac16 U+\frac56 D,\frac12 L+\frac12 M\right)
\]

\[
\text{For }a=5: \ (D,L),\ (U,R),\ \left(\frac16 U+\frac56 D,\frac12 L+\frac12 M\right),\ \left(pU+(1-p)D,\ R\right),\ p\in\left[\frac23,1\right]
\]
